\documentclass[12pt]{article}
\usepackage[a4paper, margin=1in]{geometry}
\usepackage{graphicx} % Required for inserting images
\usepackage[backend=bibtex]{biblatex}
\usepackage{hyperref}
\usepackage{dirtree}
\title{\textbf{Mini Game Hub}}
\author{Larissa \\ Sahil Deo}
\date{\today}
\addbibresource{references.bib}
\begin{document}
\maketitle
\section{External Libraries}
\begin{itemize}
\item \textbf{numpy} - Used for efficient handling of arrays and game board matrix. It simplifies operations on the game board, such as checking valid moves, updating states, and checking the win condition.

\item \textbf{pygame} - Used for building the graphical user interface of the games. It handles rendering, mouse detection, game loops, and animations.

\item \textbf{matplotlib} - Used for visualizing game statistics such as top 5 winners and the frequencies of each game

\item \textbf{time} - Used for implementing delays, timers, and tracking time-based events within games.

\item \textbf{os} - Used for interacting with the file system via bash (os.system())

\item \textbf{sys} - Used for taking command line arguments for game.py 
\end{itemize}

\section{Design \& Architecture}

\subsection{Authentication System (Bash)}
User authentication is handled via shell scripts:
\begin{itemize}
    \item User credentials are stored in \texttt{.users.tsv} with SHA-256 hashed passwords.
    \item Login allows 3 attempts before temporary lockout.
    \item Lockout timestamps are stored in \texttt{.failed\_logins.csv}.
\end{itemize}

\subsection{Leaderboard System}
Game statistics are stored in CSV files. The leaderboard script:
\begin{itemize}
    \item Sorts users by wins, losses, or win/loss ratio
    \item Formats output using \texttt{column}
    \item Maintains a top-5 players file
\end{itemize}


\subsection{Game Class Abstraction}
The \texttt{Game} class stores:
\begin{itemize}
    \item The board as a NumPy 2D array
    \item Current player turn (boolean)
    \item Board dimensions and cell size
\end{itemize}
It provides utility functions such as \texttt{switchTurn()}, \texttt{valid\_move()}, and a pluggable \texttt{checkWin()} function passed during initialization, allowing reuse across different games.

\subsection{Rendering and Scaling}
All rendering is done using \texttt{pygame}. Custom scaling functions \texttt{scale\_h()} and \texttt{scale\_w()} ensures resolution-independent scaling by mapping coordinates relative to a reference screen size.

Static elements such as boards are pre-rendered onto \texttt{Surface} objects and blitted every frame to improve performance.

\subsection{Event Handling and Game Loop}
Each game runs a main loop that:
\begin{itemize}
    \item Captures user input (mouse clicks, key presses)
    \item Updates game state (moves, turn switching, timers)
    \item Renders updated visuals
\end{itemize}
Mouse clicks are mapped to board cells, and moves are executed only if valid.

\subsection{Win Detection}
Each game defines its own \texttt{checkWin()}:
\begin{itemize}
    \item Tic Tac Toe and Connect 4 scan in 4 directions (horizontal, vertical, both diagonals) using NumPy slicing.
    \item Othello determines the winner by counting pieces when no valid moves remain (using \texttt{np.sum()})
\end{itemize}

\subsection{Animations}
Animations are time-based using timestamps:
\begin{itemize}
    \item Tic Tac Toe: progressive drawing of crosses (line segments) and circles (arc sweep).
    \item Connect 4: falling discs simulated using $y = \frac{1}{2}gt^2$.
    \item Othello: flipping discs using \texttt{pygame.draw.ellipse()} and expansion.
\end{itemize}

Jagged drawing effects are implemented via the \texttt{JaggedLine} class, which perturbs line segments randomly for a hand-drawn aesthetic.
\texttt{pygame.draw.arc()} is naturally jagged, which is what inspired the hand-drawn aesthetic.

\subsection{Blitz Mode}
A \texttt{Timer} class manages turn-based countdowns. Each player’s timer activates only during their turn. The remaining time is displayed as a dynamically coloured vertical bar, transitioning from green to red as time decreases.

\subsection{User Interface Components}
Reusable UI elements include:
\begin{itemize}
    \item \texttt{Button}: handles hover, click feedback, and rendering
    \item \texttt{Box}: used for popups and text displays
\end{itemize}
End-of-game overlays use semi-transparent surfaces to dim the screen and display options like ``Play Again'' and ``Main Menu''.

\subsection{Othello Move Logic}
Valid moves are computed by scanning in all 8 directions from each empty cell. A move is valid if it captures at least one opponent piece bracketed by the current player's pieces. Piece flipping is performed by updating slices of the NumPy array.

\subsection{Connect 4 Mechanics}
Each column tracks the number of filled cells. When a move is made, the disc is placed in the lowest available row. Falling animation is applied until the disc reaches its final position.

\subsection{Flowchart:}

\begin{center}
\begin{tabular}{c}
Start \\
$\downarrow$ \\
Authentication (main.sh) \\
$\downarrow$ \\
Select Game (game.py)\\
$\downarrow$ \\
Choose Time Control (game.py)\\
$\downarrow$ \\
Toss to Decide First Player (game.py)\\
$\downarrow$ \\
\begin{tabular}{ccc}
TicTacToe & Connect4 & Othello \\
play() & play() & play()
\end{tabular} \\
$\downarrow$ \\
Return Result \\
$\downarrow$ \\
Leaderboard (leaderboard.sh) \\
$\downarrow$ \\
End
\end{tabular}
\end{center}

\section{Features Implemented}
\begin{itemize}
    \item \textbf{Authentication System} :
    \begin{itemize}
        \item Supports user registration and login with passwords securely hashed using SHA-256.
        \item Two-player authentication is required before starting a game.
        \item The system enforces a 3-attempt login limit per user, after which a temporary time-based lockout is applied.
    \end{itemize}
    
    \item \textbf{Blitz Feature} :
    \begin{itemize}
        \item A timed gameplay mode where each player must make a move within a specified time limit.
        \item The remaining time is visually represented using a vertical timer bar which changes colour according to the urgency of move.
    \end{itemize}


    \item \textbf{Animations} : Visual enhancements implemented to improve user experience:
    \begin{itemize}
        \item Tic Tac Toe: drawing animations for X and O symbols, along with a line animation highlighting the winning combination.
        \item Othello: coin flipping with a wave-like animation to represent piece conversion visually.
        \item Connect Four: a falling animation for coins to simulate realistic gravity-based placement.
    \end{itemize}

    \item \textbf{Abstractions} : 
    \begin{itemize}
        
        \item \textbf{Developer Point of View}
            \begin{itemize}
                \item Design elements which are inheritable as a Class Button, Box, Jagged Line etc.
                \item Menus - callable as a function, returns the choice                 
            \end{itemize}
        \item \textbf{User Point of View}
            \begin{itemize}
                \item GameHub folder is clean - system files and user files are hidden
                \item Intuitive Menus and usage             
            \end{itemize}
        
    \end{itemize}
    

\end{itemize}


\section{Directory Structure}

\dirtree{%
.1 \texttt{hub/}.
.2 \texttt{main.sh}.
.2 \texttt{game.py}.
.2 \texttt{leaderboard.sh}.
.2 \texttt{reset.sh}.
.2 \texttt{games/}.
.3 \texttt{tictactoe.py}.
.3 \texttt{connect4.py}.
.3 \texttt{othello.py}.
.3 \texttt{classGame.py}.
.3 \texttt{design\_elements.py}.
.2 \texttt{images/}.
.3 \texttt{red\_disc.png}.
.3 \texttt{yellow\_disc.png}.
.2 \texttt{.user\_files/}.
.3 \texttt{.stats\_tictactoe.csv}.
.3 \texttt{.stats\_connect4.csv}.
.3 \texttt{.stats\_othello.csv}.
.3 \texttt{.users.tsv}.
.3 \texttt{.user\_total\_wins.csv}.
.3 \texttt{.top5.txt}.
.3 \texttt{.game\_frequencies.csv}.
.2 \texttt{.system\_files/}.
.3 \texttt{.failed\_logins.csv}.
}

\section{Installation \& Setup}
\begin{verbatim}
pip install numpy pygame matplotlib
\end{verbatim}

\section{Usage}
\texttt{bash main.sh}

\section{Retrospective}
\begin{itemize}
    \item Accidental change of the same part by both teammates - \\
    \emph{Fix:} Clearly communicating beforehand, and also doing a merge commit when necessary
    \item Lagging animation - \\
    \emph{Fix:} Performance improved by precomputing static frames onto Surface objects and then blitting them to the screen at loop runtime \cite{pygame_surface}
    \item Timer and Toss Bar fluid going out of box for low levels - \\
    \emph{Fix:} Using pygame Surfaces transparency feature - made a window Surface \\ which lets viewer peek through only the shape we want, hence clipping the fluid \cite{pygame_transparency}
    \item Scaling of sizes across systems - \\
    \emph{Fix:} Using scaling functions wherever necessary
    \item Not able to change game board values in Othello - the diagonal array extracted from game board using the diagonal() function was not writable by default.\\
    \emph{Fix:} Used arr\_name.setflags(write=1) \cite{numpy_setflags}
    \item Seed was same for all board lines since time was converted to int, and execution was finished in less than a second. \\
    \emph{Fix:} Used int(time\_ns())
    \item Unnecessary complicated scaling by introducing a Size class - \\
    \emph{Fix:} Removed overcomplicated Size class, changed it to simple scaling functions 

\end{itemize}

\section{What we would do differently with 10 more hours}

\begin{itemize}

\item \textbf{Code Architecture Improvements}
\begin{itemize}
\item Introduce a state-machine type flow for game.py to simplify the logic 
\item Create a Menu Class with Button and Box as subclass
\item Use enums for things like whose turn is it to avoid all the confusion during development
\end{itemize}

\item \textbf{User Experience Enhancements}
\begin{itemize}
\item Add smoother scene transitions (fade in/out between menus and games).
\item Improve aesthetics and feedback (hover effects, button animations, sound effects).
\item Add a pause button
\end{itemize}

\item \textbf{Feature Expansion}
\begin{itemize}
\item Add single-player mode which just chooses random choices (Just for fun)
\item Figure out a cleaner leaderboard display aesthetic so we can see terminal output simultaneously
\end{itemize}

\end{itemize}

\nocite{*}
\printbibliography
\end{document}
